\section{Architecture and Threat Model}
Fig.~\ref{fig:architecture} shows the proposed architecture. The system places MNO policy at the SMO/Non-RT RIC layer and uses bounded agents for O-RAN security automation. The telemetry agent observes network and media-plane indicators. The reasoning agent maps state to trust level and candidate action. The policy agent enforces the action ladder. The proof agent supports bounded admissibility checks. The audit/anchor agent writes evidence hashes and compact commitments.

\begin{figure*}[!t]
\centering
\begin{tikzpicture}[
font=\small,
box/.style={draw, rounded corners, thick, align=center, minimum height=8mm, minimum width=24mm},
arrow/.style={-{Latex[length=2mm]}, thick},
node distance=8mm
]
\node[box, fill=trustblue!18] (mno) {MNO / Tenant\\Policy Intent};
\node[box, fill=trustblue!18, right=of mno] (smo) {SMO / Non-RT RIC\\Policy Registry};
\node[box, fill=trustgreen!18, right=of smo] (nrt) {Near-RT RIC\\xApps};
\node[box, fill=trustorange!18, right=of nrt] (ran) {O-CU / O-DU / O-RU\\RTP / DTLS-RTP};

\node[box, fill=trustpurple!18, below=12mm of smo] (obs) {Telemetry\\OBSERVE};
\node[box, fill=trustpurple!18, right=of obs] (reason) {Reasoning\\REASON};
\node[box, fill=trustpurple!18, right=of reason] (policy) {Policy\\GATE};
\node[box, fill=trustpurple!18, right=of policy] (proof) {Proof\\PROVE};
\node[box, fill=trustpurple!18, right=of proof] (audit) {Audit/Anchor\\ANCHOR};

\node[box, fill=trustred!12, below=12mm of policy] (zk) {AUTH\_V2.x\\Admissibility};
\node[box, fill=trustred!12, left=of zk] (evidence) {Evidence Store\\CID / Hash};
\node[box, fill=trustred!12, right=of zk] (chain) {Audit Contract\\Commit / Reject};

\draw[arrow] (mno)--node[above]{intent}(smo);
\draw[arrow] (smo)--node[above]{A1 policy}(nrt);
\draw[arrow] (nrt)--node[above]{E2 control}(ran);
\draw[arrow] (ran.south)|-(obs.east);
\draw[arrow] (obs)--(reason);
\draw[arrow] (reason)--(policy);
\draw[arrow] (policy)--(proof);
\draw[arrow] (proof)--(audit);
\draw[arrow] (proof)--(zk);
\draw[arrow] (audit)--(evidence);
\draw[arrow] (audit)--(chain);
\draw[arrow] (chain.north)|-(nrt.south);
\end{tikzpicture}
\caption{Proposed ZKTrustLLM-Agents L4 architecture. Policy is placed in the SMO/Non-RT RIC governance path while bounded agents connect O-RAN telemetry to proof, evidence, and audit anchoring.}
\label{fig:architecture}
\end{figure*}

\begin{table}[!t]
\centering
\caption{Bounded Agentic Components}
\label{tab:agents}
\begin{adjustbox}{width=\columnwidth}
\begin{tabular}{lll}
\toprule
Agent & Role & Output \\
\midrule
Telemetry & Observe O-RAN/media KPIs & Jitter, loss, profile \\
Reasoning & Map state to trust/action & Trust state, action class \\
Policy & Enforce action ladder & Auto/Human/Privileged/Never \\
Proof & Check admissibility & AUTH\_V2.x status \\
Audit/Anchor & Bind evidence and commit & CID/hash/anchor record \\
\bottomrule
\end{tabular}
\end{adjustbox}
\end{table}


\subsection{Threat Controls}
The framework considers unapproved autonomous action, evidence tampering, anchor replay, zero/empty anchors, unsafe escalation, and unauthorized submitters. Fig.~\ref{fig:threats} maps these concepts to controls.

\begin{figure}[!t]
\centering
\begin{tikzpicture}[
font=\scriptsize,
phase/.style={draw, rounded corners, thick, fill=trustgreen!18, align=center, minimum width=15mm},
threat/.style={draw, rounded corners, thick, fill=trustred!15, align=center, minimum width=20mm},
arrow/.style={-{Latex[length=1.6mm]}, thick},
node distance=4mm
]
\node[phase] (o) {OBSERVE};
\node[phase, right=of o] (r) {REASON};
\node[phase, right=of r] (p) {PROVE};
\node[phase, right=of p] (a) {ANCHOR};
\node[phase, right=of a] (act) {ACT};
\draw[arrow] (o)--(r);
\draw[arrow] (r)--(p);
\draw[arrow] (p)--(a);
\draw[arrow] (a)--(act);
\node[threat, below=8mm of r] (t1) {Unsafe\\action};
\node[threat, below=8mm of p] (t2) {Invalid\\proof};
\node[threat, below=8mm of a] (t3) {Replay /\\zero anchor};
\draw[arrow, trustred] (t1)--node[left]{policy gate}(r);
\draw[arrow, trustred] (t2)--node[right]{AUTH check}(p);
\draw[arrow, trustred] (t3)--node[right]{contract reject}(a);
\end{tikzpicture}
\caption{Threat-to-control mapping. Unsafe actions are policy-gated, invalid proofs are blocked, and replay/zero anchors are rejected.}
\label{fig:threats}
\end{figure}


\subsection{Smart-Contract Boundary}
The smart contract does not control O-RAN traffic directly. Its role is trust-plane accountability: storing compact commitments, rejecting duplicate commitments, rejecting zero commitments, and enforcing authorized submitters. It cannot verify physical network truth, guarantee evidence availability, or replace MNO operational policy.
